\documentclass[uplatex,a4paper,12pt]{jsarticle}
\usepackage{amsmath,amssymb}
\usepackage{enumitem}

\begin{document}

\title{線形代数 確認テスト 解答}
\author{}
\date{}
\maketitle

\section*{配点}
全10問，各10点，合計100点．60点以上を合格とする．

\section*{解答}

\begin{enumerate}

\item
正解は
\[
\boxed{(b)}
\]
である．$(1,2,3)^T$ はベクトルである．

\item
\[
\|\mathbf{x}\|_2
=
\sqrt{3^2+4^2}
=
\sqrt{9+16}
=
5
\]
したがって，
\[
\boxed{5}
\]

\item
\[
\mathbf{a}\cdot\mathbf{b}
=
1\cdot4+2\cdot0+3\cdot2
=
4+0+6
=
10
\]
したがって，
\[
\boxed{10}
\]

\item
\[
A+B=
\begin{pmatrix}
1+5 & 2+6\\
3+7 & 4+8
\end{pmatrix}
=
\begin{pmatrix}
6&8\\
10&12
\end{pmatrix}
\]
したがって，
\[
\boxed{
\begin{pmatrix}
6&8\\
10&12
\end{pmatrix}
}
\]

\item
\[
A^T=
\begin{pmatrix}
1&4\\
2&5\\
3&6
\end{pmatrix}
\]
したがって，
\[
\boxed{
\begin{pmatrix}
1&4\\
2&5\\
3&6
\end{pmatrix}
}
\]

\item
\[
AB=
\begin{pmatrix}
1&2\\
3&4
\end{pmatrix}
\begin{pmatrix}
2&0\\
1&3
\end{pmatrix}
\]
\[
=
\begin{pmatrix}
1\cdot2+2\cdot1 & 1\cdot0+2\cdot3\\
3\cdot2+4\cdot1 & 3\cdot0+4\cdot3
\end{pmatrix}
=
\begin{pmatrix}
4&6\\
10&12
\end{pmatrix}
\]
したがって，
\[
\boxed{
\begin{pmatrix}
4&6\\
10&12
\end{pmatrix}
}
\]

\item
\[
A\circ B=
\begin{pmatrix}
1\cdot2 & 2\cdot0\\
3\cdot1 & 4\cdot3
\end{pmatrix}
=
\begin{pmatrix}
2&0\\
3&12
\end{pmatrix}
\]
したがって，
\[
\boxed{
\begin{pmatrix}
2&0\\
3&12
\end{pmatrix}
}
\]

\item
正解は
\[
\boxed{(b),(c)}
\]
である．

行列積は一般に交換法則 $AB=BA$ を満たさない．
一方で，結合法則
\[
A(BC)=(AB)C
\]
と分配法則
\[
A(B+C)=AB+AC
\]
は成り立つ．

\item
\[
A=
\begin{pmatrix}
1&2\\
2&4
\end{pmatrix}
\]
では，2行目が1行目の2倍である．

したがって，独立な行は1本なので，
\[
\boxed{\operatorname{rank}(A)=1}
\]

\item
正解は
\[
\boxed{(a)}
\]
である．

固有値 $\lambda$，固有ベクトル $p$ は，
\[
Ap=\lambda p
\]
を満たす．

\end{enumerate}

\end{document}